\section{Statements}\label{sec:results}

\subsection{The packing specification}\label{sec:results-general}

\begin{definition}[Packing specification]\label{def:intro-spec}
\lean{SymplecticPacking.Problem.PackingSpec}
Let $n\ge1$, $k\ge0$, $C>0$, $c=(c_1,\dots,c_k)\in(0,\infty)^k$ and $s=(s_1,\dots,s_k)\in(\R^{2n})^k$. A
function $H\colon\R\times\R^{2n}\to\R$ satisfies $\Pack(n,k,C,c,s)$ if
\begin{enumerate}[label=(\alph*)]
\item $H$ is $C^\infty$, and there is a compact set $\mathcal S\subset\R^{2n}$ with $H(t,z)=0$ for all
 $t\in\R$ and all $z\notin\mathcal S$;
\item for every $i$ and every $z\in B(s_i,c_i)$ there is an integral curve $\gamma\colon\R\to\R^{2n}$ of
 $X_H$ with $\gamma(0)=z$, and every integral curve $\gamma\colon[0,1]\to\R^{2n}$ of $X_H$ with
 $\gamma(0)=z$ satisfies $\gamma(1)\in\Bo(C)$.
\end{enumerate}
\leanmark{Problem.PackingSpec}
\end{definition}

For $H$ satisfying (a) the flow $\phi^t_H$ is global and consists of symplectic diffeomorphisms, and
(b) says $\phi^1_H(B(s_i,c_i))\subset\Bo(C)$ for every $i$.

\begin{definition}[Admissible data]\label{def:intro-admissible}
\lean{SymplecticPacking.Problem.Admissible}
Let $n\ge2$, $k\ge1$, $C>0$ and $c\in(0,\infty)^k$. The data $(n,k,C,c)$ are \emph{admissible} if
\begin{align*}
 n\ge3:&\qquad \textstyle\sum_ic_i^n<C^n\quad\text{and}\quad c_i+c_j<C\ \text{ for all } i\ne j;\\
 n=2:&\qquad \textstyle\sum_ic_i^2<C^2\quad\text{and}\quad 3c_i<C\ \text{ for all } i.
\end{align*}
\leanmark{Problem.Admissible}
\end{definition}

For $n=2$ the condition $3c_i<C$ comes from the fibre capacity $\frac13$ of the model; it is a sufficient
small-capacity range, not the full four-dimensional packing criterion. This paper uses admissibility only for $n=2$
and equal capacities.

\subsection{Robustness}\label{sec:intro-contribution}

Theorem~\ref{thm:intro-4d} leaves an explicit margin: the images lie in the closed ball of radius $R$, while the
density requirement is already met by a larger radius $R'$. The following estimate converts this margin into a
tolerance on the vector field.

\begin{mainproposition}[Robustness]\label{thm:intro-stability}
Let $k\ge10$, $\eps\in(0,1)$, $R=(k/(1-\eps/4))^{1/4}$ and $R'=(k/(1-\eps/2))^{1/4}$, so $R<R'$. Let $\widehat H$ be a smooth function on $\R\times\R^4$
vanishing outside $\R\times K$ for a compact set $K$. Suppose that, for $t\in[0,1]$, the field $X_{\widehat H}(t,\cdot)$
is $L$-Lipschitz for some $L\ge0$ and $|X_{\widehat H}-X_{H_{k,\eps}}|\le\eta$ on $\R^4$, where $\eta e^L\le R'-R$. Then the time-one map
of $\widehat H$ sends each closed unit ball $B(s_i,\pi)$ into the closed ball of radius $R'$ about the origin, and
$k/R'^4=1-\eps/2>1-\eps$.
\leanmark{four\_dim\_stability}
\end{mainproposition}

\begin{proof}
Let $\gamma$ and $\hat\gamma$ be integral curves of $X_{H_{k,\eps}}$ and $X_{\widehat H}$ starting at the same point of
$B(s_i,\pi)$. Then $|\hat\gamma(1)-\gamma(1)|\le\eta(e^L-1)/L\le\eta e^L$ by Gr\"onwall's inequality
(Lemma~\ref{lem:eff-gronwall}, with the Lipschitz bound on $X_{\widehat H}$), and $|\gamma(1)|\le R$ by
Theorem~\ref{thm:intro-4d}. For $L=0$ the Gr\"onwall bound is $\eta$.
\end{proof}

No explicit bound on the Lipschitz constant of $X_{H_{k,\eps}}$ is given, so the margin $R'-R$ is not
directly a numerical step-size bound.

\begin{lemma}[Gr\"onwall]\label{lem:eff-gronwall}
Let $G,H\colon\R\times\R^{2n}\to\R$ be $C^\infty$, let $\eta>0$ and $L_X>0$, and suppose that
\begin{enumerate}[label=(\roman*)]
\item $|X_G(t,z)-X_H(t,z)|\le\eta$ for all $t\in[0,1]$ and all $z\in\R^{2n}$;
\item $X_H(t,\cdot)$ is $L_X$-Lipschitz on all of $\R^{2n}$ for every $t\in[0,1]$.
\end{enumerate}
If $\gamma_G,\gamma_H\colon[0,1]\to\R^{2n}$ are integral curves of $X_G$ and $X_H$ with
$\gamma_G(0)=\gamma_H(0)$, then $|\gamma_G(1)-\gamma_H(1)|\le\eta(e^{L_X}-1)/L_X$.
\end{lemma}

\begin{proof}
$u=\gamma_G-\gamma_H$ satisfies $u(0)=0$ and $|\dot u|\le\eta+L_X|u|$, so Gr\"onwall's inequality gives the bound.
\end{proof}

For smooth compactly supported $H$, (ii) holds for every $L_X>0$ with $L_X\ge\max_{[0,1]\times\R^{2n}}\|D^2_zH\|$, as
$X_H=J_0\nabla_zH$ with $J_0$ orthogonal.

\subsection{The Epoch data}\label{sec:results-epoch}

\begin{maincorollary}[The Epoch data]\label{thm:main-epoch}
Let $k\ge10$ be an integer, $\eps\in(0,\frac78)$, $s_i=(3i-3,0,0,0)\in\R^4$, $R=(k/(1-\eps/2))^{1/4}$ and $C=\pi R^2$.
Then $H=H[2,k,C,(\pi,\dots,\pi),s]$ satisfies $\Pack(2,k,C,(\pi,\dots,\pi),s)$: its time-one map sends each closed unit
ball $B(s_i,\pi)$ into $\Bo(\pi R^2)$, and $k/R^4=1-\eps/2>1-\eps$.
\leanmark{epoch\_packing, epoch\_theorem\_ten}
\end{maincorollary}

Theorem~\ref{thm:intro-4d} is Corollary~\ref{thm:main-epoch} at $\eps/2$. The bound $\eps<\frac78$ is the range
$\rho\sqrt k=\sqrt{1-\eps/2}>\frac34$ of the linear comparison datum. The companion paper~\cite{SongHD} proves the
corollary for all $k\ge9$ and $\eps\in(0,1)$, with the scalar comparison datum outside this range, and uses it as the
input of its finite-output algorithm.

For $H\colon\R\times\R^4\to\R$, $t\in\R$ and $h>0$, write $\mathrm{RK4}(H,t,h,y)=y+\frac h6(k_1+2k_2+2k_3+k_4)$ for
one classical fourth-order Runge--Kutta step of $X_H$, where $k_1=X_H(t,y)$, $k_2=X_H(t+\frac h2,y+\frac h2k_1)$,
$k_3=X_H(t+\frac h2,y+\frac h2k_2)$, and $k_4=X_H(t+h,y+hk_3)$. The radial margin $R'-R>0$ between
Corollary~\ref{thm:main-epoch} and Theorem~\ref{thm:intro-4d} yields uniform eventual containment for perturbed
fixed-step RK4 applied directly to $H_{k,\eps}$.

\begin{maincorollary}[Uniform perturbed fixed-step RK4 containment]\label{cor:rk4-containment}
Let $k\ge10$ be an integer and $\eps\in(0,1)$, and put $H=H_{k,\eps}$, $R_Q=(k/(1-\eps/4))^{1/4}$,
$R=(k/(1-\eps/2))^{1/4}$, and $g=R-R_Q>0$. There exist an integer $N_0\ge1$ and a real $\delta_0>0$ such that for every
integer $N\ge N_0$, every $1\le i\le k$, every source point $x\in B(s_i,\pi)$, and every sequence
$y_0,y_1,\dots,y_N\in\R^4$ satisfying
\[
 |y_0-x|\le\delta_0 \qquad\text{and}\qquad
 \bigl|y_{m+1}-\mathrm{RK4}\bigl(H,\tfrac mN,\tfrac1N,y_m\bigr)\bigr|\le\frac{\delta_0}{N}
 \quad(0\le m<N),
\]
the terminal iterate satisfies $|y_N|\le R$, with $k/R^4=1-\eps/2>1-\eps$. Explicitly, for any $K,C\ge0$ such that
$|X_H(t,x)-X_H(s,x)|\le K|t-s|$, $|X_H(t,x)-X_H(t,y)|\le K|x-y|$, and $|X_H(t,y)|\le C$ for all $s,t\in[0,1]$ and
$x,y\in\R^4$, one may take $M=K+KC$ and
\[
 \delta_0 = \frac{g}{4e^K},\qquad N_0 = \left\lceil\frac{4Me^K}{g}\right\rceil + 1.
\]
\leanmark{epoch\_rk4\_convergence, epoch\_exact\_rk4\_convergence}
\end{maincorollary}

\begin{proof}
Since $H\in C^\infty(\R\times\R^4)$ has compact spatial support uniformly in $t$, finite global bounds $K,C\ge0$ exist
on $[0,1]\times\R^4$. Each sub-stage $k_j$ ($j=2,3,4$) evaluates $X_H$ within temporal distance $\frac1N$ and spatial
distance $\frac CN$ of $(\frac mN,y_m)$, so $|\frac16(k_1+2k_2+2k_3+k_4)-X_H(\frac mN,y_m)|\le\frac{5M}{6N}\le\frac MN$.
Combined with the per-step state perturbation $\le\delta_0/N$, the effective step velocity $v_m=N(y_{m+1}-y_m)$ is
within $\delta_0+M/N$ of $X_H(\frac mN,y_m)$, and the discrete Gr\"onwall grid estimate along the exact integral curve
$\gamma$ from $x$ ($|\gamma(1)|\le R_Q$ by Theorem~\ref{thm:intro-4d}) gives
$|y_N-\gamma(1)|\le(|y_0-x|+\delta_0+2M/N)e^K\le 2\delta_0e^K+\frac{2Me^K}{N}\le\frac g2+\frac g2=g$, whence
$|y_N|\le R_Q+g=R$.
\end{proof}

Here $K$ and $C$ are obtained from compactness and are not numerically instantiated, and a first-order consistency
estimate suffices for eventual containment without invoking fourth-order smoothness constants; accordingly the proof does
not provide executable step or precision budgets. It establishes target containment; certification of an implemented
\texttt{float64} Hamiltonian and of the grader's nearby-pair numerical symplecticity check requires additional bounds.

\subsection{Explicitness}\label{sec:intro-explicitness}

\begin{definition}[Definite description]\label{def:intro-the}
\lean{SymplecticPacking.Explicit.the}
For a property $\mathcal P$ of elements of a set $X$ and a proof $h$ that exactly one $x\in X$ satisfies
$\mathcal P(x)$, $\defdesc(h)$ denotes that element. Phrases such as ``the unique $u$ with $F(u)=y$''
denote this operation.
\end{definition}

\begin{definition}[Explicit definitional construction]\label{def:intro-levels}
An \emph{explicit definitional construction} (level L1) is a definition built from rational constants and the input by
arithmetic and elementary functions, integrals over compact domains, limits of series with an explicit rate of
convergence, maxima and minima of continuous functions over explicit compact sets, unique solutions of explicit
ordinary differential equations, inverses of explicit injective maps, unique roots of explicit equations, case
distinctions on propositions about the input and about the arguments of the function being defined, and least
elements of searches over a fixed enumeration of finitely described objects (integers, rationals, and tuples of these
and of rational polynomials) whose validity is decided by finitely many explicit inequalities; definite descriptions
are used only after unique existence has been proved (Definition~\ref{def:intro-the}).
\end{definition}

Theorem~\ref{thm:intro-4d} concerns level L1 (Table~\ref{tab:ham-explicit}). The finite form and computability
(levels L2 and L3) are defined in~\cite{SongHD}.
